\section{Dataset}
\label{sec:dataset}

\subsection{Frameworks}

We ingest nine AI security and governance frameworks spanning risk
management, technical controls, threat modeling, and compliance guidance:

\begin{enumerate}
    \item \textbf{NIST AI RMF}~\citep{nist_ai_rmf} --- 76 controls
          covering AI risk governance and lifecycle management.
    \item \textbf{MITRE ATLAS}~\citep{mitre_atlas} --- 218 entries
          cataloging adversarial tactics, techniques, and procedures
          targeting AI systems.
    \item \textbf{OWASP Top~10 for LLM Applications}~\citep{owasp_llm_top10}
          --- 10 high-level vulnerability categories for large language
          model deployments.
    \item \textbf{OWASP Agentic AI Security Risks} --- 10 risk
          categories specific to autonomous AI agent architectures.
    \item \textbf{OWASP AI Exchange} --- 88 entries providing a
          cross-cutting taxonomy of AI security threats and mitigations.
    \item \textbf{CSA AI Controls Matrix}~\citep{csa_aicm} --- 261
          controls organized by domain (IAM, data security, model
          governance, etc.) for cloud-hosted AI workloads.
    \item \textbf{CoSAI Risk Model} --- 61 risk entries from the
          Coalition for Secure AI covering supply chain, inference,
          and deployment risks.
    \item \textbf{AI Use Case Framework (AIUC)} --- 189 entries
          cataloging governance requirements across AI application
          categories.
    \item \textbf{EU GPAI Code of Practice} --- 70 entries derived
          from the European Union's General-Purpose AI Code of Practice
          for providers of foundation models.
\end{enumerate}

Together these yield 983 control nodes. Framework sizes range from 10
entries (OWASP LLM, OWASP Agentic) to 261 (CSA AICM), reflecting the
heterogeneous granularity of real-world AI governance standards.

\begin{figure}[t]
\centering
\includegraphics[width=\linewidth]{fig_categorical_composition}
\caption{Categorical composition of the corpus. (A)~Entry types per
framework showing absolute counts. (B)~Overall entry-type distribution.
(C)~Edge confidence distribution across the knowledge graph.}
\label{fig:categorical_composition}
\end{figure}

\subsection{Graph Structure and Data Sources}

The control knowledge graph integrates edges from four complementary
sources, each capturing a different type of inter-control relationship:

\begin{enumerate}
    \item \textbf{Within-framework hierarchies and cross-framework
          graph edges} (5,813 edges total: 1,090 intra-framework
          parent-child relationships plus 4,723 cross-framework edges
          derived from high-confidence BM25 bridge matching).
    \item \textbf{Upstream community mappings} from the GenAI Security
          Project crosswalk~\citep{genai_security_project} (3,210
          total rows; 209 in-scope after filtering to the nine
          included frameworks), used under CC~BY-SA~4.0.
    \item \textbf{Upstream cross-reference tables} (363 total rows;
          58 in-scope), capturing additional control relationships
          published alongside framework documentation.
    \item \textbf{Bootstrap-pruned anchor pairs} from 12 pair-config
          files (816 total anchor pairs), generated via co-citation
          analysis followed by bootstrap cross-validation pruning to
          remove unstable candidates.
\end{enumerate}

After deduplication, these four sources produce cross-framework edges
spanning 15 unique framework pairs out of the $\binom{9}{2} = 36$
possible pairs, yielding a graph connectivity of 41.7\%.
Figure~\ref{fig:cross_framework_heatmap} visualizes the direction-agnostic
edge distribution across all framework pairs.

\begin{figure}[t]
\centering
\includegraphics[width=\linewidth]{fig_cross_framework_heatmap}
\caption{Cross-framework edge counts (direction-agnostic). Left: heatmap
of pairwise edge counts after symmetrizing. Right: marginal totals per
framework. CSA AICM and AIUC dominate connectivity due to their large
control inventories and extensive upstream mappings.}
\label{fig:cross_framework_heatmap}
\end{figure}

\begin{figure}[t]
\centering
\includegraphics[width=\linewidth]{fig_coverage_completeness}
\caption{Cross-framework coverage completeness. Each cell shows the
fraction of possible control pairs between two frameworks that have at
least one edge in the unified knowledge graph. Coverage ranges from
dense (e.g., AIUC--CSA AICM) to zero for unconnected pairs.}
\label{fig:coverage_completeness}
\end{figure}

\subsection{Labeling Protocol}

We employ a two-stage labeling protocol:
\begin{enumerate}
    \item \textbf{LLM-SME silver labels.} A structured prompt elicits
          tier judgments from a large language model acting as a
          simulated subject matter expert. Each pair receives a tier
          label (0--3) and a free-text rationale. 1,515 training pairs
          and 379 validation pairs are labeled this way.
    \item \textbf{Expert gold labels.} A human expert independently
          labels 400 pairs using the same four-tier scheme. These form
          the frozen test set (\texttt{v1\_frozen}) used for all
          reported metrics. The frozen split is SHA-verified and never
          used during training or hyperparameter selection.
\end{enumerate}

\subsection{Tier Scheme}

The four tiers capture decreasing degrees of semantic overlap:
\begin{description}
    \item[Equivalent (3):] Controls address the same requirement with
          comparable specificity.
    \item[Related (2):] Controls address overlapping concerns but differ
          in scope or granularity.
    \item[Partial (1):] One control partially addresses the other's
          requirement.
    \item[Unrelated (0):] No meaningful semantic overlap.
\end{description}

The frozen test set is class-imbalanced: Related pairs dominate
(186/400 = 46.5\%), with Unrelated and Partial each at 78/400 (19.5\%)
and Equivalent at 58/400 (14.5\%).
